\begin{figure}[t!]
	\caption{\textbf{Anti-regime civil society organization activity and democratic mobilizations}}
	\footnotesize
	\vspace{-.25cm}
	\renewcommand{\baselinestretch}{1}
    \captionsetup{singlelinecheck=off}
    \captionsetup{width=1\linewidth}
	\caption*{This figure presents an event study comparing majority Catholic autocracies to non-Catholic autocracies in their anti-regime civil society organization (CSO) activity and frequency of democratic mobilizations and protests as determined by indices from the V-Dem database.  The reference year is set to 1959, the first year of the doctrinal shift. Endpoints are binned and are not shown. The anti-regime CSO activity index ranks the threat posed by anti-regime civil society organizations on a scale of 0 to 4, where 0 is no anti-regime civil society organization activity, and 4 is a major present threat to the governing regime from anti-regime civil society organizations. The democratic mobilization index assesses the number of small- and large-scale demonstrations in favor of democracy in a given year with a maximum value of 4. The autocracy designation is also constructed from V-Dem data, and includes all closed or electoral autocracies from their ``regimes of the world'' variable. Data on the percentage of the population that is Catholic comes from the World Religion Project. These data are extended backward using the first year of data. The vertical grey bars show the treatment window from 1959--1963. Country and year fixed effects are included. The red bars represent a 90\% confidence interval with standard errors clustered by country.}
	\label{fig:treatmentEffect}
	\footnotesize
	\vspace{-.5cm}
    \begin{center}
	\begin{tabular}{cc}
	\textbf{Panel~A: Anti-regime CSO activity} & \textbf{Panel~B: Democratic mobilizations}  \\
	\includegraphics[width=.45\textwidth]{../figures/raw/figure_5a_event_study_anti_system_cso.pdf}& \includegraphics[width=.45\textwidth]{../figures/raw/figure_5b_event_study_democratic_protests.pdf}
	\end{tabular}
	\end{center}
\end{figure}
